\section{Chapitre~\ref{sec:dofa}. Apprendre avec \nt{DOPA}}

Les mécanismes de la synapse (quanta, détecteur de coïncidences, calcium, fenêtres de plasticité) et le comportement de la dopamine comme erreur de prédiction sont mesurés directement; que les règles mènent au gradient et ce que coûte la diffusion sont des théorèmes; le résultat de l'apprentissage du choix est un calcul par le modèle du livre; le schéma qui calcule l'erreur est une hypothèse.

{\footnotesize
\setlength{\tabcolsep}{3pt}\renewcommand{\arraystretch}{1.15}
\begin{longtable}{>{\raggedright}p{0.5cm}>{\raggedright}p{4.0cm}>{\raggedright}p{5.6cm}>{\raggedright}p{2.3cm}>{\raggedright\arraybackslash}p{1.7cm}}
\toprule
N\textsuperscript{o} & Affirmation & Expérience et ce qui est mesuré & Source & Statut \\
\midrule\endfirsthead
\toprule
N\textsuperscript{o} & Affirmation & Expérience et ce qui est mesuré & Source & Statut \\
\midrule\endhead
1 & La rétropropagation de l'erreur, sous la forme qu'elle a dans les réseaux artificiels, n'a pas été trouvée dans le cerveau & Pas d'expérience directe: c'est une affirmation d'absence. Des revues analysent ce qu'exige l'algorithme (connexions de retour reproduisant les connexions aller, phase séparée) et les approximations biologiques proposées. & \cite{Lillicrap2020, WhittingtonBogacz2019} & hypothèse \\
2 & Le poids se décompose en nombre de sites de libération, probabilité de libération et réponse à un quantum: $w=N_s\,p\,A_1$~\eqref{eq:quantal} & Jonction neuromusculaire de grenouille à libération réduite: la réponse du destinataire varie par paliers multiples de la réponse miniature spontanée à un quantum; le nombre de quanta par impulsion suit une loi de Poisson. & \cite{delCastillo1954} & mesuré \\
3 & Le récepteur du destinataire est un détecteur de coïncidences; le calcium déclenche la modification du poids & Patch-clamp de neurones de souris: les ions Mg$^{2+}$ bouchent au potentiel de repos le canal ouvert par le glutamate et le libèrent lors de la dépolarisation. Tranches d'hippocampe: un chélateur du calcium dans le destinataire bloque la potentialisation à long terme, et la libération photo-induite de calcium dans la cellule renforce à elle seule la transmission. & \cite{Nowak1984, Malenka1988calcium} & mesuré \\
4 & N'importe quel côté exécute la décision: $A_1$ croît chez le destinataire, $p$ change chez l'émetteur & Le glutamate libéré par la lumière au-dessus d'une seule épine provoque une croissance rapide et durable de l'épine et du courant de ses récepteurs; la potentialisation s'accompagne de l'insertion de récepteurs AMPA. La dépolarisation du destinataire supprime temporairement la libération chez l'émetteur; l'effet est porté par des endocannabinoïdes qui retraversent la fente (montré sur des entrées \nt{GABA}). La dépression à court terme dépend de la probabilité de libération de l'émetteur. & \cite{Matsuzaki2004, Malinow2002, WilsonNicoll2001, Tsodyks1997} & mesuré \\
5 & La synapse se renforce quand l'émetteur et le destinataire sont actifs ensemble~\eqref{eq:hebb} & Lapin anesthésié: après de courtes séries d'impulsions sur la voie d'entrée de l'hippocampe, la réponse est renforcée pendant $30$~min à $10$~h. En tranche d'hippocampe, les impulsions du destinataire ne renforcent que les synapses actives au même moment. & \cite{Bliss1973LTP, Kelso1986hebb} & mesuré \\
6 & La règle~\eqref{eq:oja} trouve le vecteur propre principal des corrélations des entrées (proposition~\ref{prop:oja}) & Théorème; démonstration dans le chapitre. Aucune expérience ne montre un neurone ayant appris la composante principale de ses entrées; le mécanisme qui limite la croissance des poids dans la synapse n'est pas établi. & \cite{Oja1982} & théorie \\
7 & Hebb temporel: fenêtre d'environ $\pm20\ms$ (fig.~\ref{fig:stdp}); des séquences répétées font pousser des chaînes & Paires de neurones d'hippocampe en culture: une impulsion du destinataire dans les $20\ms$ qui suivent celle de l'émetteur donne un renforcement durable, dans les $20\ms$ qui précèdent, un affaiblissement; les deux dépendent des récepteurs NMDA. Le rapport $A_-=0.6A_+$ de la figure est une illustration. Que des chaînes poussent ainsi dans un réseau n'a pas été mesuré directement. & \cite{BiPoo1998} & mesuré \\
8 & Trace~\eqref{eq:threefactor}: la dopamine arrivée $1$--$2\,\text{s}$ après l'activité conjointe renforce la connexion, plus tard non (proposition~\ref{prop:threefactor}) & Tranches de striatum, stimulation optogénétique séparée des entrées glutamatergiques et dopaminergiques: l'épine ne grossit que si la dopamine arrive $0.3$--$2\,\text{s}$ après l'entrée glutamatergique; la fenêtre est fixée par la montée rapide et la dégradation de l'AMPc. Dans le cortex, l'appariement laisse des traces séparées pour le renforcement et l'affaiblissement, que les récepteurs des monoamines transforment en modifications durables dans une fenêtre de l'ordre de la seconde. & \cite{Yagishita2014, He2015eligibility} & mesuré \\
9 & Des fenêtres de plusieurs secondes existent aussi sans dopamine: le troisième signal est une forte excitation du destinataire & Souris vivante, champ CA1: un seul événement de plateau dans le destinataire renforce les entrées arrivées quelques secondes avant et après lui, et crée un champ de lieu en un seul passage. & \cite{Bittner2017} & mesuré \\
10 & Le pas moyen de la règle à trois facteurs suit le gradient de la récompense~\eqref{eq:perturbation} (proposition~\ref{prop:gradient}) & Théorème du gradient pour l'apprentissage par écarts aléatoires; dans le chapitre, établi pour un bruit faible. & \cite{Williams1992} & théorie \\
11 & Le bruit est une recherche: le réseau apprend de ses propres écarts aléatoires & Jeunes diamants mandarins: l'inactivation du noyau de sortie du circuit des ganglions de la base supprime brusquement et réversiblement la variabilité du chant. Diamants adultes: une perturbation est appliquée aux exécutions d'une syllabe dont la hauteur est d'un côté de la moyenne, et les oiseaux déplacent vite la hauteur de l'autre côté: ils apprennent d'après leur propre dispersion. & \cite{Olveczky2005, TumerBrainard2007} & mesuré \\
12 & Le choix entre deux options s'apprend avec $\tau_e=1\,\text{s}$ et $\delta=R-\bar R$; pas avec $\tau_e=50\ms$; sans soustraction, les poids croissent sans limite, et avec une grande vitesse d'apprentissage le réseau peut figer le plus mauvais choix & \texttt{models/dofa\_learning.py}, $\eta=0.05$, $500$ essais, $6$ graines. $\tau_e=1\,\text{s}$: part des choix de $A$ $0.65$--$0.79\to0.96$--$1.00$, poids $0.85$--$1.33$. $\tau_e=50\ms$: $0.38$--$0.55$, poids inchangés. $\delta=R$: poids du vainqueur $860$--$1190$. $\delta=R$, $\eta=0.5$, $3$ graines: l'une s'est figée sur $B$ ($0.04\to0$). Le script imprime les quatre cas; le recalcul concorde avec le chapitre. & \texttt{models/\allowbreak dofa\_\allowbreak learning.py} & modèle \\
13 & Le temps d'apprentissage par un seul signal commun croît proportionnellement au nombre de neurones bruités (proposition~\ref{prop:broadcastcost}) & Courbes d'apprentissage d'un réseau linéaire: la perturbation des neurones est plus lente que le gradient exact d'un facteur égal au nombre de neurones de sortie, la perturbation des poids encore d'un facteur égal au nombre d'entrées. & \cite{Werfel2005} & théorie \\
14 & Les sorties de \nt{DOPA} vont avant tout aux régions qui choisissent les actions & Reconstruction complète des axones de neurones \nt{DOPA} nigrostriés isolés chez le rat: les branches d'une cellule couvrent $0.45$--$5.7\,\%$ du volume du striatum (en moyenne $2.7\,\%$), et il n'y a presque pas de branches hors du striatum. & \cite{Matsuda2009} & mesuré \\
15 & Le cortex apprend à prédire son entrée, l'erreur est calculée sur place & Revue: le cortex sensoriel présente des réponses au désaccord entre l'entrée attendue et l'entrée reçue. Qu'il s'agisse d'une règle générale d'apprentissage du cortex n'est pas démontré. & \cite{Keller2018} & hypothèse \\
16 & La dopamine se comporte comme l'erreur~\eqref{eq:ctd}: pic, transfert sur le signal annonciateur, creux (fig.~\ref{fig:ctdcases}) & Neurones \nt{DOPA} du singe: pic pour un jus inattendu; après apprentissage, pic sur le signal annonciateur et pas de réponse au jus promis; creux quand le jus n'arrive pas. La forme continue~\eqref{eq:ctd} relève de la théorie. & \cite{Schultz1997, Doya2000} & mesuré \\
17 & Schéma qui calcule $\dd V/\dd t$ et soustrait l'attente & Souris, enregistrement et optogénétique dans l'aire tegmentale ventrale: les neurones \nt{DOPA} soustraient l'attente de la récompense; les neurones \nt{GABA} voisins les inhibent quand la récompense est attendue, et les exciter ou les inhiber modifie l'erreur. Le calcul de la dérivée n'a pas été montré. & \cite{Eshel2015arithmetic} & hypothèse \\
18 & Le neurone \nt{DOPA} est un pacemaker; les creux sont moins profonds que les pics & En tranche, les neurones \nt{DOPA} déchargent spontanément et très régulièrement, sur un potentiel pacemaker lent. Chez le singe, la fréquence suit quantitativement l'erreur positive, mais ne transmet que faiblement l'erreur négative. & \cite{GraceOnn1989, Bayer2005} & mesuré \\
19 & L'acteur et le critique (fig.~\ref{fig:actorcritic}) & L'algorithme vient de la théorie de l'apprentissage par renforcement. Chez l'homme (IRMf), le striatum ventral se comporte comme un critique avec ou sans choix, le striatum dorsal seulement quand il faut choisir des actions. & \cite{SuttonBarto2018, ODoherty2004actor} & hypothèse \\
20 & Le signe de $\delta$ dans la règle est inversé chez les neurones de la seconde famille de récepteurs & Tranches de striatum: chez les neurones à récepteurs D1, l'appariement produit un renforcement qui exige D1; chez les neurones à D2, un affaiblissement qui exige D2. & \cite{Shen2008} & mesuré \\
\bottomrule
\end{longtable}}
